\section{怎样掌握文言通假字}

中学生在学习文言文时，经常碰到一些字形熟悉却既不能用本义，也不能用引申义来解释的字，老师就说这个字是另外某个字的通假字。例如贾谊《论积贮疏》中有这样一段话：

“兵旱相乘，天下大屈，有勇力者聚徒而衡击，晏夫赢老易子而咬其骨；政治未毕通也，远方之能疑者，并举而争起矣。”

这段话的意思是说：兵灾、旱灾交互侵袭，社会上（财富）非常缺乏，有胆量有力气的人就聚众抢劫，体弱年老的人就把孩子交换着吃。政治还没有完全上轨道，离朝廷远的地方那些僭立为王的人，就一齐争着起来造反了。

这段话中不好讲的，是那三个加点的字。老师会说：“衡”通“横”，“罢”通“疲”，“疑”通“拟”（“疑”的简体字）。同学们就会问：怎么知道这个字通那个字呢？

是的，所谓“通假”并不是随心所欲地用这个字代替那个字，其中是有规律可寻的。我们只要能掌握通假字的规律，再在阅读文言文时根据上下文意融会理解，就能判断出通假字。为了说明的方便，也考虑实际的需要，我们把通假字范围划大一些，把“古今字”（指古代原有的字和后出的字）也包括在内。下面我们就说说通假字的规律。

1. 某些“声旁字”可以代替包含这个“声旁字”的“形声字”。形声字有形旁和声旁两部分，在古书中往往把一些形声字用它们的声旁字来代替。例如：

\noindent
\begin{tabular}{llll}
	汝——女， & 值——直， & 返——反， & 舍——舍，\\
	俸——奉， & 蓄——畜， & 拟——疑， & 胸——匈，\\
	座——坐， & 材——才， & 诚——戒， & 庭——廷，\\
	娶——取， & 俱——具， & 属——属， & 彰——章，\\
	否——不， & 悬——县， & 擒——禽， & 授——受，\\
	播——雷， & 曝——暴， & 智——知， & 裂——列，\\
	履——厌， & 缮——善， & 纵——从， & 影——景，\\
	债（价）——贯， & 岸——干， & 溢——益，& 捧——奉，\\
	侯——宾， & 缠——虚， & 愬——忽， & 酬（划）——画，\\
	洁——絮， & 徊——回， & 熟——敦，& 燃——然，\\
	锲——契， & 钜——巨， & 陌——百，& 暮——莫，\\
	剂——齐， & 性——生， & 烫——汤，& 纹——文，\\
	响——向， & 债——责， & 肢——支，& 猝——卒，\\
	伙——火， & 现——见， & 镝——质，& 噱——转……\\
\end{tabular}

如此等等，不一而足。这就是说，古人常用一个“声旁字”来代替包含这个“声旁字”的“形声字”，现代人就把这种现象叫做“通假”，即“声旁字”在意义上与包含这个“声旁字”的“形声字”相通。如说“女”通“汝”，“直”通“值”，“反”通“返”，等等。“通”的意思是说，这个字的意义与那个字的意义通用。在古书注释中，或在字典、词典的义项中，常用“×通×”，这是最明确的对于通假字的说法；有时也写成“×读日×”、“×读为×”、“×读如×”、“×读若×”、“×与×同”。

我们从上面所说的“声旁字”代替“形声字”的那些例字中，可以看出：大多数情况下，把古代的“声旁字”加个偏旁部首就成为今天大家熟悉的字，语言学家就把那些古代的声旁字叫做“古字”，而把加了偏旁部首的后起字叫做“今字”。对于初学古代汉语的中学生和一般社会青年来说，不必细分哪些是古今字，哪些是同音代替的假借字，一律用通假字这个统称，掌握起来更方便些。我们在学习文言文时，如果发现某个字在句子中不论从其本义或引申义来讲都讲不通，就可以从它加了偏旁部首的同音字考虑，也许上下文意就豁然贯通了。

2. 某些“声旁”相同的“形声字”可以在意义上相同。一个形声字，既然包含声旁和形旁两部分，那么，某些声旁相同而形旁不同的形声字在意义上也就可以互相替代，这也是通假字中常见的一种现象。例如：
\begin{enumerate}[label=\textcircled{\arabic*}]
\item “陇”通“垄”（高地、山冈）；无陇断焉。（《愚公移山》）
\item “厝”通“措”（放置）；一厝朔东，一厝雍南。（《愚公移山》）
\item “帖”通“贴”（粘）；对镜帖花黄。（《木兰诗》）
\item “油”通“屈”（弯曲）；卧右膝，油右臂支船。（《核舟记》）
\item “板”通“版”（用木板制成供印刷用的底子）；有布衣毕升，又为活板。（《活板》）
\item “径”通“经”（经过）；江水又东径黄牛山下。（《三峡》）
\item “娉”通“聘”（旧时以礼物订婚）；云是当为河伯妇，即娉取。（《西门豹治邺》）
\item “徧”通“遍”（遍及）；小惠未徧，民弗从也。（《曹刿论战》）
\item “谄”通“悦”（高兴）；秦王不说。（《唐睢不辱使命》）
\item “唱”通“倡”（倡导）；为天下唱，宜多应者。（《陈涉世家》）
\item “适”通“谪”（强迫调发）；发间左适戍渔阳九百人。（陈涉世家》）
\item “被”通“披”（穿着）；将军身被坚执锐。（《陈涉世家》）
\item “距”通“拒”（抵挡，拒守）；子墨子九距之。（《公输》）

距关，毋内诸侯。（《鸿门宴》）

\item “辨”通“辩”（辩解）；不复——自辨。（《答司马谏议书》）
\item “辩”通“辨”（分辨）；万钟，则不辩礼义而受之。（《鱼我所欲也》）
\item “踰”通“逾”（越过）；喻淮之精也……（《甘渚疏序》）
\item “逝”通“誓”（发誓）；逝将去女，适彼乐土。（《硕鼠》）
\item “案”通“按”（审察，查看）；召有司案图。（《廉颇蔺相如列传》）

 “案”通“按”（按住）；熟卧空室之中，若有所畏惧，则梦见夫人据案其身哭矣。（《订鬼》）
 
\item “輮”通“煣”（使······弯曲）；不复挺者，輮使之然也。（《劝学》）
\item “栖”通“杯”（杯子）；沛公不胜桎梏，不能辞。（《鸿门宴》）
\item “扣”通“叩”（敲打）：得双石于潭上，扣而聆之。（《石钟山记》）
\item “喩”通“险”（险要）：予观雁荡诸峰，皆峭拔峻怪。（《雁荡山》）
\item “蠟”通“缣”（细密）：古之治天下，至蠟至悉也。（《论积贮疏》）
\item “殴”通“驱”（驱使）：今殴民而归之农。（《论积贮疏》）
\item “详”通“倖”（假装）：乃令张仪详去秦。（《屈原列传》）
\item “椎”通“锤”（锤子）：若见鬼把椎锁绳缰，立守其旁。（《订鬼》）
\item “缙”通“摴”（插）：缙绅、大夫、士萃于左丞相府。（《指南录后序》）
\item “陵”通“凌”（凌侮）：几为巡徼所陵迫死。（《指南录后序》）
\item “钞”通“抄”（抄写）：道中手自钞录。（《指南录后序》）
\item “霾”通“埋”（陷没）：霾两轮兮繁四马。（《国殇》）
\item “衙”通“炫”（炫耀）：将衙外以惑愚瞽也？（《卖柑者言》）
\item “凛凛”通“凛凛”（危惧的样子）：可以为富安天下，而直为此凛凛也！（《论积贮疏》）
\item “呐”通“讷”（说话迟钝）：我将按酒入去，只听得差拨口里响出一句“高太尉”三个字来。（《林教头风雪山神庙》）
\item “决”通“诀”（诀别）：行过夷门，见侯生，具告所以欲死秦军状。辞诀而行。（《信陵君窃符救赵》）
\end{enumerate}

3.\ 某些形体不同而读音相同或相近的字，在意义上可以相通。通假字的根本原则是在用字上同音替代。这有点像现代人写别字似地用一个同音字去替代另一个字，不过古人的这种偶尔遗忘造成的同音替代现象，我们叫它“通假字”，而且一旦有前人这样用，后人也就可以跟着用，这就与现代人写别字完全不同了。我们掌握了这个原则之后，就可以把文言文中讲不通的某个字与它的同音字联系起来，并且根据上下文考虑哪一个同音字的意义恰当，那么，就可能确定那个讲不通的某个字的通假字。例如：

\begin{enumerate}[label=\textcircled{\arabic*}]
\item “惠”通“慧”（聪明）；甚矣，汝之不惠。（《愚公移山》）
\item “胥”通“须”（等待）；左师触龙愿见太后，太后盛气而胥之。（《触龙说赵太后》）
\item “麾”通“挥”（指挥）；麾众拥豪民马前。（《书博鸡者事》）
\item “止”通“只”（只，不过）；虎因喜，计之曰：“技止此耳！”（《黔之驴》）
\item “尔”通“耳”（罢了）；我亦无他，惟手熟尔。（《卖油翁》）
\item “食”通“饲”（喂）；食马者，不知其能千里而食也。（《马说》）
\item “邪”通“耶”（吗）；呜呼！其真无马邪？其真不知马也！（《马说》）
\item “少”通“稍”（稍微）：寨中人蜷伏不少动。（《冯婉贞》）
\item “帅”通“率”（率领）：帅突将三千为前驱。（《李愬雪夜入蔡州》）
\item “衡”通“横”（纵横的横）：内立法度，务耕织，修守战之具，外连衡而斗诸侯。（《过秦论》）
\item “已”通“以”（表示时间、方位、数量的界线）：自董卓已来，豪杰并起。（《隆中对》）
\item “以”通“已”（已经）：卒买鱼烹食，得鱼腹中书，固以怪之矣。（《陈涉世家》）
\item “阙”通“缺”（中断）：自三峡七百里中，两岸连山，略无阙处。（《三峡》）
\item “圉”通“御”（防御）：子墨子之守圉有余。（《公输》）
\item “熙”通“嬉”（戏耍）：圣人非所与熙也。（《晏子使楚》）
\item “炎”通“焰”（火焰）：顷之，烟炎张天。（《赤壁之战》）
\item “有”通“又”（表示重复和继续）：虽有槁暴，不复挺者，使之然也。（《劝学》）
\item “要”通“邀”（邀请）：张良出，要项伯。（《鸿门宴》）
\item “倍”通“背”（背弃）：愿伯具言臣之不敢倍德也。（《鸿门宴》）
\item “蚤”通“早”（早些）：旦日不可不蚤自来谢项王。（《鸿门宴》）
\item “邻”通“隙”（隔阂）；今者有小人之言，令将军与臣有隙。（《鸿门宴》）
\item “匍”通“奏”（大声）；洞天石扉，匍然中开。（《梦游天姥吟留别》）
\item “悦”通“恍”（恍惚）；悦惊起而长嗟。（《梦游天姥吟留别》）
\item “擎”通“揽”（持，握）；两峰秀色，俱可手擎。（《游黄山记》）
\item “而”通“尔”（你的）；而翁归，自与汝复算耳！（《促织》）
\item “裁”通“才”（刚刚）；“手栽举，则又超忽而跃”。（《促织》）
\item “翼”通“翌”（明天）；翼日进宰，宰见其小，怒呵成。（《促织》）
\item “侥”通“俯”（低下）；百越之君，侥首系颈，委命下吏。”（《过秦论》）
\item “鸠”通“纠”（集合）；鸠宗族僮仆百许人。（《书博鸡者事》）
\item “得”通“德”（感恩）；今为所识穷乏者得我而为之。（《鱼我所欲也》）
\item “矢”通“屎”（粪便）；矢溺皆闭其中，与饮食之气相搏。（《狱中杂记》）
\item “函胡”通“含糊”（不清晰）；南声函胡，北音清越。（《石钟山记》）
\item “乡”通“向”（向着）；请数公子行日，以至晋鄙军之日北乡自到，以送公子！（《信陵君窃符救赵》）
\item “拔”通“攀”（拉着）；父利其然也，日扳仲永环谒于邑人，不使学。（《伤仲永》）
\item “简”通“拣”（选择）；盖简桃核修狭者为之。（《核舟记》）
\item “卤”通“鲁”（莽撞）；重念蒙君实视遇厚，于反复不宜卤莽。（《答司马谏议书》）
\item “顿”通“钝”（磨损）；必以全争于天下，故兵不顿而利可全。（《谋攻》）
\item “竦”通“耸”（高高地直立）；水何澹澹，山岛竦峙。（《观沧海》）
\item “父”通“甫”（古代常用作人的表字）；四人者：庐陵萧君圭君玉，长乐王回深父，余弟安国平父、安上纯父。（《游褒禅山记》）
\item “别”通“瘞”（物体表面凹下去）；颧骨望上翘，嘴唇往下别。（《孙悟空三打白骨精》）
\item “阳”通“伴”（假装）；皆阳应曰：“诺。”（《记王忠肃公翱事》）
\end{enumerate}

除了以上列举的音同通假的例字之外，还有一些上古读音相同或具有双声、叠韵的关系，而在今天读音不同的异形字，也可以互相通假。这当然给我们辨析通假字带来一定困难，不过这一类通假字为数不多，我们只要下功夫记一记就可以解决了。例如：

\begin{enumerate}[label=\textcircled{\arabic*}]
\item “还”通“旋”（回转）；扁鹊望桓侯而还走。（《扁鹊见蔡桓公》）
\item “趣”通“促”（催促）；巫妪何久也？弟子趣之。（《西门豹治邺》）
\item  “识”通“志”（记住）：因笑谓迈曰：“汝识之乎？······”（《石钟山记》）
\item  “内”通“纳”（纳入）：距关，毋内诸侯。（《鸿门宴》）
\item  “信”通“仲”（伸张）：孤不度德量力，欲信大义于天下。”（《隆中对》）
\item  “罢”通“疲”（疲乏）：罢夫赢老易子而咬其骨。（《论积贮疏》）
\item  “猗”通“兮”（语气助词）：坎坎伐檀兮，置之河之干兮，河水清且涟猗。（《伐檀》）
\item  “盖”通“盍”（何，怎么）：技盖至此乎？（《庖丁解牛》）
\item  “亡”通“无”（没有）：河曲智叟亡以应。（《愚公移山》）
\item  “没”通“昧”（冒着）：愿令得补黑衣之数，以卫王宫，没死以闻。（《触龙说赵太后》）
\item  “生”通“性”（禀性、资质）：君子生非异也，善假于物也。（《劝学》）
\item  “景”通“影”（影子）：天下云集响应，赢粮而景从。（《过秦论》）
\end{enumerate}

诚然，我们学习文言文，通假字常常是一个难点。但是，只要我们掌握了一些有关通假字的知识，懂得通假字的基本规律，再根据上下文意仔细琢磨，就能判断出通假字。
\newpage